\medskip
\noindent
14. \underline{Links with transcendental theory: the rigid-analytic
case.  The canonical homomorphism}
$\underline M_K\longrightarrow A_K^{\natural\mathrm o}$.

\medskip
\noindent
\SGARefTarget{sga7.sga7i.eix.n14-1}{14.1}. We assume here the foundations of the theory of rigid-analytic
spaces, for which one may consult [12], [13], and [34].

Let $S$ be a complete trait, and let $A_K$ be an abelian scheme over its
field of fractions.  In \SGARefLink{sga7.sga7i.eix.n8-1}{8.1}, by descent from the Raynaud extension in
the semistable case, we constructed an extension
$G_K=A_K^{\natural\mathrm o}$ (\SGARefLink{sga7.sga7i.eix.n8-1-1}{8.1.1})

\begin{equation}\tag{\SGARefTarget{sga7.sga7i.eix.n14-1-1}{14.1.1}}
 0\longrightarrow T_K\longrightarrow G_K\longrightarrow B_K
 \longrightarrow0
\end{equation}

\noindent
of an abelian scheme $B_K$ by a torus $T_K$.  Considering the dual
abelian scheme $A'_K$, we likewise have an extension
$G'_K=A'{}^{\natural\mathrm o}_K$

\begin{equation}\tag{\SGARefTarget{sga7.sga7i.eix.n14-1-2}{14.1.2}}
 0\longrightarrow T'_K\longrightarrow G'_K\longrightarrow B'_K
 \longrightarrow0\qquad ,
\end{equation}

\noindent
where $B'_K$ is moreover identified with the abelian scheme dual to
$B_K$ (\SGARefLink{sga7.sga7i.eix.n8-2}{8.2}).  As in (\SGARefLink{sga7.sga7i.eix.n8-1-8}{8.1.8}), let $\underline M_K$ and
$\underline M'_K$ denote the lattices over $K$ dual to $T'_K$ and
$T_K$, respectively:

\begin{equation}\tag{\SGARefTarget{sga7.sga7i.eix.n14-1-3}{14.1.3}}
 \underline M_K=D(T'_K),\qquad \underline M'_K=D(T_K)\qquad .
\end{equation}

\noindent
Thus the extension structures (\SGARefLink{sga7.sga7i.eix.n14-1-2}{14.1.2}) and (\SGARefLink{sga7.sga7i.eix.n14-1-1}{14.1.1}) may respectively be
interpreted (VIII \SGARefLink{sga7.sga7i.eviii.n3-7}{3.7}) as equivalent to the data of homomorphisms

\begin{equation}\tag{\SGARefTarget{sga7.sga7i.eix.n14-1-4}{14.1.4}}
 \underline M_K\xrightarrow{\ p_K\ }B_K,\qquad
 \underline M'_K\xrightarrow{\ p'_K\ }B'_K\qquad .
\end{equation}

Let $W_K^{\natural}$ be the biextension of $(B_K,B'_K)$ by
$\mathbb G_{mK}$ which defines the canonical duality between these
abelian schemes (\SGARefLink{sga7.sga7i.eix.n8-2}{8.2}), and let $E_K$ be the biextension of
$(\underline M_K,\underline M'_K)$ by $\mathbb G_{mK}$ obtained by
pulling $W_K^{\natural}$ back along the homomorphisms (\SGARefLink{sga7.sga7i.eix.n14-1-4}{14.1.4}).  One
checks immediately that trivializations of $E_K$ correspond
canonically and bijectively to homomorphisms

\begin{equation}\tag{\SGARefTarget{sga7.sga7i.eix.n14-1-5}{14.1.5}}
 i:\underline M_K\longrightarrow G_K
\end{equation}

\noindent
lifting $p_K:\underline M_K\longrightarrow B_K$, or, symmetrically, to
homomorphisms

\begin{equation}\tag{\SGARefTarget{sga7.sga7i.eix.n14-1-6}{14.1.6}}
 i':\underline M'_K\longrightarrow G'_K
\end{equation}

\noindent
lifting $p'_K:\underline M'_K\longrightarrow B'_K$.

With this understood, write a superscript ${}^{\mathrm{an}}$ for the
rigid-analytic groups over $K$ associated with the algebraic groups
under consideration [13].  In [24] one defines a canonical exact
sequence

\begin{equation}\tag{\SGARefTarget{sga7.sga7i.eix.n14-1-7}{14.1.7}}
 0\longrightarrow\underline M_K^{\mathrm{an}}
 \xrightarrow{\ i^{\mathrm{an}}\ }G_K^{\mathrm{an}}
 \xrightarrow{\ e\ }A_K^{\mathrm{an}}\longrightarrow0\qquad .
\end{equation}

\noindent
Here $i^{\mathrm{an}}$, which may also be interpreted as a homomorphism
$i:\underline M_K\longrightarrow G_K$ of ordinary group schemes,
\underline{lifts the homomorphism} $p$ in (\SGARefLink{sga7.sga7i.eix.n14-1-4}{14.1.4}), that is, it is of
the form considered in (\SGARefLink{sga7.sga7i.eix.n14-1-5}{14.1.5}).  This exact sequence is the
rigid-analytic analogue of the exact sequence (\SGARefLink{sga7.sga7i.eix.n13-3-4}{13.3.4}).  It depends
functorially on $A_K$, and its formation is compatible with every
change of traits $S'\longrightarrow S$.

The latter fact also shows that, in order to define and study (\SGARefLink{sga7.sga7i.eix.n14-1-7}{14.1.7}),
finite descent reduces us to the case where $A_K$ has semistable
reduction.  In this case the group schemes
$T_K,G_K,B_K,\underline M_K,\ldots$ are the generic fibres of group
schemes $T,G,B,M,\ldots$ over $S$, with $G$ an extension of the abelian
scheme $B$ by the torus $T$ (paragraph 7).  We may then consider the
formal completion $\hat G$ of $G$ along its special fibre, and by
construction there is a canonical isomorphism (\SGARefLink{sga7.sga7i.eix.n7-2-6}{7.2.6})

\begin{equation}\tag{\SGARefTarget{sga7.sga7i.eix.n14-1-8}{14.1.8}}
 \hat G\xrightarrow{\ \sim\ }\widehat{A^{\mathrm o}}
 \qquad .
\end{equation}

\noindent
Passing to the ``generic fibres'', which are rigid-analytic groups over
$K$, gives a canonical isomorphism

\begin{equation}\tag{\SGARefTarget{sga7.sga7i.eix.n14-1-9}{14.1.9}}
 \hat G^{\mathrm{an}}\xrightarrow{\ \sim\ }
 \widehat{A^{\mathrm o}}\,{}^{\mathrm{an}}
 \qquad .
\end{equation}

\noindent
On the other hand, there are canonical homomorphisms (rigid-analytic
open immersions, but not in general isomorphisms)

\begin{equation}\tag{\SGARefTarget{sga7.sga7i.eix.n14-1-10}{14.1.10}}
 \hat G^{\mathrm{an}}\lhook\joinrel\longrightarrow G_K^{\mathrm{an}},
 \qquad
 \widehat{A^{\mathrm o}}\,{}^{\mathrm{an}}
 \lhook\joinrel\longrightarrow A_K^{\mathrm{an}}
 \qquad .
\end{equation}

\noindent
The homomorphism $e$ in (\SGARefLink{sga7.sga7i.eix.n14-1-7}{14.1.7}) is uniquely characterized by the
commutativity of the diagram

\begin{equation}\tag{\SGARefTarget{sga7.sga7i.eix.n14-1-11}{14.1.11}}
\begin{tikzcd}
 \hat G^{\mathrm{an}} \arrow[r] \arrow[d]
   & \widehat{A^{\mathrm o}}\,{}^{\mathrm{an}} \arrow[d] \\
 G_K^{\mathrm{an}} \arrow[r] & A_K^{\mathrm{an}}
\end{tikzcd}
\end{equation}

\noindent
where the unnamed arrows are those of (\SGARefLink{sga7.sga7i.eix.n14-1-9}{14.1.9}) and (\SGARefLink{sga7.sga7i.eix.n14-1-10}{14.1.10}).

\medskip
\noindent
\underline{Remarks \SGARefTarget{sga7.sga7i.eix.n14-1-12}{14.1.12}.}

\smallskip
\noindent
a) Notice that the uniqueness of $e$ making the preceding diagram
commute is immediate, since the rigid-analytic space
$G_K^{\mathrm{an}}$ is connected.  The existence of $e$, on the other
hand, is by no means trivial, and on this point we merely refer to the
cited note of M.~Raynaud.  Moreover, the definition of
$i^{\mathrm{an}}$ given in loc. cit. is likewise transcendental in
nature, although it may be interpreted as an object of an essentially
algebraic nature (\SGARefLink{sga7.sga7i.eix.n14-1-5}{14.1.5}), once the construction of $G$ itself, and
hence of $G_K$, has been obtained through the constructions of
\underline{formal geometry} in paragraph 7.  I know of no purely
algebro-geometric description of it, but below (\SGARefLink{sga7.sga7i.eix.n14-4}{14.4}) I shall propose
an ``arithmetic'' characterization, obtained by reducing to a situation
of finite type over the ring $\mathbb Z$.  Thus, from the standpoint of
algebraic geometry, the homomorphism (\SGARefLink{sga7.sga7i.eix.n14-1-5}{14.1.5}) seems to be a more hidden
piece of structure than those encountered in the preceding paragraphs.

\smallskip
\noindent
b) As explained above, the homomorphism $i$ is canonically associated
with a well-determined trivialization of the biextension $E_K$ of
$(\underline M_K,\underline M'_K)$, or equivalently with a
well-determined homomorphism (\SGARefLink{sga7.sga7i.eix.n14-1-6}{14.1.6}) lifting $p'_K$.  One then proves
that the latter is precisely the homomorphism occurring in the exact
sequence

\begin{equation}\tag{\SGARefTarget{sga7.sga7i.eix.n14-1-13}{14.1.13}}
 0\longrightarrow\underline M'{}^{\mathrm{an}}_K
 \xrightarrow{\ i'{}^{\mathrm{an}}\ }G'{}^{\mathrm{an}}_K
 \longrightarrow A'{}^{\mathrm{an}}_K\longrightarrow0
\end{equation}

\noindent
analogous to (\SGARefLink{sga7.sga7i.eix.n14-1-7}{14.1.7}), associated with the abelian scheme $A'_K$ dual
to $A_K$ (unless it is the opposite???).
